\chapter{Pipelines, predicates and dispatch}\label{ch:2}
\section*{The working problem}
Build small reusable transformations with explicit inputs and outputs. Operator syntax is helpful only when the underlying function contracts remain visible.

Pipes carry values forward. Predicates combine Boolean decisions. Match tables select handlers by keys or patterns. Although all three use templates, their execution and fallback policies are different.

\section*{Concrete callables and ownership}
The pipe and predicate wrappers retain callable objects. Capturing a value makes the callable independent of the local variable, while capturing a reference couples its validity to another object's lifetime. A composed object may be returned from a factory, stored in a member or used after its defining scope has exited. Review captures at that ownership boundary, not only at the line containing the lambda.

\section*{In this chapter}
\begin{enumerate}
\item Ordered pipe stages
\item Predicate conjunction
\item Negation and alternatives
\item Keyed match tables
\item Whole-input pattern checks
\end{enumerate}
Read the workshops in order, then use the chapter review to connect the small examples to an application boundary. Each example is independently runnable; it does not depend on variables introduced in an earlier listing.

\clearpage
\workshop{2.1}{Ordered pipe stages}{C++20}
A pipeline is a left-to-right data dependency. Each stage accepts the preceding result, so changing stage order changes semantics even if every stage has the same host type. Keep stages small enough that a failing intermediate value is easy to inspect.

\lstinputlisting[language=C++,caption={Example ex02\_01: C++}]{examples/cpp/ex02_01.inc}
\section*{Execution, step by step}
Two becomes five and then twenty. Parentheses around the complete pipeline make the assertion unambiguous to the host parser.

\section*{Read the implementation}
The relevant source entry is PipeStage in vendor/DSLtk.hpp, around line 6064. Follow the public call into its result construction or dispatch branch. Identify the input ownership, the success condition and the exact failure path before changing this example.

The assertions are the acceptance contract for this workshop. Run the complete program, then change one input at a time. When a change fails, first identify whether the failure belongs to the native build, structural recognition, language checking, execution or the domain result.

\clearpage
\workshop{2.1}{Ordered pipe stages}{D}
The D example uses a typed callable parameter. A generic callable can instead be supplied as an alias template argument. Preserve capture ownership when moving a local stage into long-lived application storage.

\lstinputlisting[language=D,caption={Example ex02\_01: D}]{examples/d/ex02_01.inc}
\section*{Failure and design boundary}
Pipes do not automatically convert exceptions into Result values. An exception thrown by a stage follows the host language rules.

\section*{Try it yourself}
Insert a stage that renders the integer to a string, then add a stage that measures its length.

\section*{Review the result}
Keep the original accepted fixture and add the new boundary fixture beside it. Compare observable values and diagnostic categories across the two hosts. Do not label the D run as passed until a D compiler has executed it; the supplied edition records source review separately from runtime evidence.

\clearpage
\workshop{2.2}{Predicate conjunction}{C++20}
A predicate wrapper turns a callable into an operand for Boolean DSL composition. Conjunction is appropriate when a value must satisfy every independent requirement, such as positivity and alignment. Pure predicates are easier to reorder and test than predicates with counters or I/O.

\lstinputlisting[language=C++,caption={Example ex02\_02: C++}]{examples/cpp/ex02_02.inc}
\section*{Execution, step by step}
Eight satisfies both tests; minus two satisfies only evenness. The resulting object remains callable with the same argument.

\section*{Read the implementation}
The relevant source entry is Predicate in vendor/DSLtk.hpp, around line 3086. Follow the public call into its result construction or dispatch branch. Identify the input ownership, the success condition and the exact failure path before changing this example.

The assertions are the acceptance contract for this workshop. Run the complete program, then change one input at a time. When a change fails, first identify whether the failure belongs to the native build, structural recognition, language checking, execution or the domain result.

\clearpage
\workshop{2.2}{Predicate conjunction}{D}
The D example uses a typed callable parameter. A generic callable can instead be supplied as an alias template argument. Preserve capture ownership when moving a local stage into long-lived application storage.

\lstinputlisting[language=D,caption={Example ex02\_02: D}]{examples/d/ex02_02.inc}
\section*{Failure and design boundary}
Do not read the ampersand as arithmetic on the argument. It combines predicate objects; their invocations perform the comparisons.

\section*{Try it yourself}
Add an upper-bound predicate and test values on both boundaries. State whether each boundary is inclusive.

\section*{Review the result}
Keep the original accepted fixture and add the new boundary fixture beside it. Compare observable values and diagnostic categories across the two hosts. Do not label the D run as passed until a D compiler has executed it; the supplied edition records source review separately from runtime evidence.

\clearpage
\workshop{2.3}{Negation and alternatives}{C++20}
Negation is a deliberate portability seam. C++ overloads logical negation for the predicate object. D uses the complement operator or the named negated method because its logical-not operator produces a Boolean rather than a new predicate.

\lstinputlisting[language=C++,caption={Example ex02\_03: C++}]{examples/cpp/ex02_03.inc}
\section*{Execution, step by step}
The outer assertion negates a Boolean result, whereas constructing nonzero negates a predicate. These operations occur at different levels.

\section*{Read the implementation}
The relevant source entry is Predicate in vendor/DSLtk.hpp, around line 3086. Follow the public call into its result construction or dispatch branch. Identify the input ownership, the success condition and the exact failure path before changing this example.

The assertions are the acceptance contract for this workshop. Run the complete program, then change one input at a time. When a change fails, first identify whether the failure belongs to the native build, structural recognition, language checking, execution or the domain result.

\clearpage
\workshop{2.3}{Negation and alternatives}{D}
The D example uses a typed callable parameter. A generic callable can instead be supplied as an alias template argument. Preserve capture ownership when moving a local stage into long-lived application storage.

\lstinputlisting[language=D,caption={Example ex02\_03: D}]{examples/d/ex02_03.inc}
\section*{Failure and design boundary}
A mechanical replacement of every exclamation mark with a tilde would corrupt ordinary Boolean tests. Port the expression according to its operand type.

\section*{Try it yourself}
Build a predicate accepting either zero or a positive even number. Test a negative odd number as a counterexample.

\section*{Review the result}
Keep the original accepted fixture and add the new boundary fixture beside it. Compare observable values and diagnostic categories across the two hosts. Do not label the D run as passed until a D compiler has executed it; the supplied edition records source review separately from runtime evidence.

\clearpage
\workshop{2.4}{Keyed match tables}{C++20}
A match table packages key-specific handlers and a fallback. Make the handlers return one coherent type so callers need not recover a variant themselves. This is a good fit for small protocol codes whose meaning is known when the program is compiled.

\lstinputlisting[language=C++,caption={Example ex02\_04: C++}]{examples/cpp/ex02_04.inc}
\section*{Execution, step by step}
The handler receives the matched key. The fallback receives an unmatched key, which allows it to construct a domain-specific error.

\section*{Read the implementation}
The relevant source entry is MatchTable in vendor/DSLtk.hpp, around line 6431. Follow the public call into its result construction or dispatch branch. Identify the input ownership, the success condition and the exact failure path before changing this example.

The assertions are the acceptance contract for this workshop. Run the complete program, then change one input at a time. When a change fails, first identify whether the failure belongs to the native build, structural recognition, language checking, execution or the domain result.

\clearpage
\workshop{2.4}{Keyed match tables}{D}
D uses module imports and exclamation-mark template instantiation where C++ uses includes, namespaces and angle brackets. The runnable file includes the required imports; the listing shows only its main body.

\lstinputlisting[language=D,caption={Example ex02\_04: D}]{examples/d/ex02_04.inc}
\section*{Failure and design boundary}
A dispatch table is not an exhaustive proof over arbitrary integers. Include a fallback unless the call boundary independently validates the key.

\section*{Try it yourself}
Add a second explicit key and replace the numeric fallback with a shared string return type.

\section*{Review the result}
Keep the original accepted fixture and add the new boundary fixture beside it. Compare observable values and diagnostic categories across the two hosts. Do not label the D run as passed until a D compiler has executed it; the supplied edition records source review separately from runtime evidence.

\clearpage
\workshop{2.5}{Whole-input pattern checks}{C++20}
Use pattern for a compact recognition predicate. The string is a compile-time argument, while the candidate input is runtime data. Whole-input recognition is suitable for field validation; prefix matching belongs in a tokenizer that explicitly advances an input position.

\lstinputlisting[language=C++,caption={Example ex02\_05: C++}]{examples/cpp/ex02_05.inc}
\section*{Execution, step by step}
The valid candidate is entirely numeric. The invalid candidate contains a legal prefix but an illegal suffix, which must not disappear during validation.

\section*{Read the implementation}
The relevant source entry is pattern in vendor/DSLtk.hpp, around line 6175. Follow the public call into its result construction or dispatch branch. Identify the input ownership, the success condition and the exact failure path before changing this example.

The assertions are the acceptance contract for this workshop. Run the complete program, then change one input at a time. When a change fails, first identify whether the failure belongs to the native build, structural recognition, language checking, execution or the domain result.

\clearpage
\workshop{2.5}{Whole-input pattern checks}{D}
D uses module imports and exclamation-mark template instantiation where C++ uses includes, namespaces and angle brackets. The runnable file includes the required imports; the listing shows only its main body.

\lstinputlisting[language=D,caption={Example ex02\_05: D}]{examples/d/ex02_05.inc}
\section*{Failure and design boundary}
The pattern facility has its own supported syntax. Do not assume every construct from a host regular-expression library is accepted.

\section*{Try it yourself}
Write three identifier boundary cases: empty input, a leading digit, and a valid name with an underscore.

\section*{Review the result}
Keep the original accepted fixture and add the new boundary fixture beside it. Compare observable values and diagnostic categories across the two hosts. Do not label the D run as passed until a D compiler has executed it; the supplied edition records source review separately from runtime evidence.

\clearpage
\section*{Chapter review: Concrete callables and ownership}
D generic lambdas can be passed as alias template parameters, for example pipe!((x) => x + x)(). Typed delegates and callable structs can be passed as ordinary values. C++ generic lambdas instead participate through their templated call operators. These are different host mechanisms. Preserve invocation behavior and lifetime first; do not demand identical source spelling. Short-circuit predicate composition also means a later predicate may not run, so side effects should not carry essential work.

\section*{Questions to answer in your implementation}
\begin{enumerate}
\item Insert a stage that renders the integer to a string, then add a stage that measures its length.
\item Add an upper-bound predicate and test values on both boundaries. State whether each boundary is inclusive.
\item Build a predicate accepting either zero or a positive even number. Test a negative odd number as a counterexample.
\item Add a second explicit key and replace the numeric fallback with a shared string return type.
\item Write three identifier boundary cases: empty input, a leading digit, and a valid name with an underscore.
\end{enumerate}
\section*{A practical completion test}
Choose one workshop and replace its fixed input with an application boundary: a command-line argument, a file, a configuration field or a typed callback. Write down the accepted input shape and the error outcome. Keep parsing and validation separate from publishing state. Preserve the smallest failing input when a regression is found.

